\subsubsection{MLP Marker}
\label{SEC:MLPMarker}
\key*{}[marker!MLP marker]

%%%%%%%%%%%%%%%%%%%%%%%%%%%%%%%%%%%%%%%%%%%%%%%%%%%%%%%%%%%%%%%%%%%%%%%%%%%%%%%%

The MLP marker is normally computed as part of the
Multidimensional Limiting Process\cite{Park_2014,Park_2016,Park_2017}.
%
It first identifies troubled cells using the
$\PO[1]$-projected MLP condition, which compares the cell vertex (corner)
values to those of the adjacent cells.
%
Then the MLP extrema detector is applied to identify local smooth extrema
near the vertex points and unmark those cells.

%%%%%%%%%%%%%%%%%%%%%%%%%%%%%%%%%%%%%%%%%%%%%%%%%%%%%%%%%%%%%%%%%%%%%%%%%%%%%%%%

When GFR outputs the MLP marker, the value \marker{} represents the
polynomial order of the limited solution,
with possible real values of [0:1,2,3,...,$\PO$].
%
Values less than $\PO$ indicate that the polynomial in the cell
has been reduced to that value to maintain stability.
%
Values between [0:1] indicate that all of the higher modes have been removed
and the linear mode was limited by that fraction in order to remain stable.

%%%%%%%%%%%%%%%%%%%%%%%%%%%%%%%%%%%%%%%%%%%%%%%%%%%%%%%%%%%%%%%%%%%%%%%%%%%%%%%%

\SeeAlso{
\inp[nmlLimiter::MarkerOption]{Marker\_Option},
\inp[nmlLimiter::MarkerVariable]{Marker\_Variable}
}


%%%%%%%%%%%%%%%%%%%%%%%%%%%%%%%%%%%%%%%%%%%%%%%%%%%%%%%%%%%%%%%%%%%%%%%%%%%%%%%%

